\documentclass[a4paper, serif, 10pt]{moderncv}

%% moderncv
% color and style
\moderncvcolor{black}
\moderncvstyle{banking}

%% page layout/margins
\usepackage[top=2.5cm, bottom=2.5cm, left=1.5cm, right=1.5cm]{geometry}

\nopagenumbers{}

%% Babel config for Brazilian Portuguese language
% ref:
% - https://latex3.github.io/babel/guides/locale-portuguese.html
\usepackage[brazilian]{babel}

%% fontspec
\usepackage{fontspec}

\IfFontExistsTF{Linux Libertine}{
  \setmainfont[Ligatures={TeX, Common}, Numbers=OldStyle]{Linux Libertine}
}{}
\IfFontExistsTF{Linux Libertine O}{
  \setmainfont[Ligatures={TeX, Common}, Numbers=OldStyle]{Linux Libertine O}
}{}

%% CTeX
% CJK support, used to show CJK characters in the resume
%
% - heading=false: disable ctex heading
% - punct=kaiming: use kaiming punctuations styles for CJK
% - scheme=plain: use plain scheme, do not override \normalsize font size
% - space=auto: space settings for CJK characters
%
% CJK fonts are resolved per locale language, e.g. Noto Serif/Sans CJK JP
% for Japanese resumes, with fallback to the next region if not installed
%
% ref:
% - http://ctan.mirrorcatalogs.com/language/chinese/ctex/ctex.pdf
\usepackage[UTF8, heading=false, punct=kaiming, scheme=plain, space=auto]{ctex}

\IfFontExistsTF{Noto Serif CJK SC}{
  \setCJKmainfont{Noto Serif CJK SC}
}{}
\IfFontExistsTF{Noto Sans CJK SC}{
  \setCJKsansfont{Noto Sans CJK SC}
}{}

%% line spacing
\usepackage{setspace}
\setstretch{1.125}

%% URL styling - use normal text instead of monospace
\urlstyle{same}

%% Patch \httplink and \httpslink to handle full URLs
% The original moderncv macros always prepend http:// or https:// to URLs.
% This causes issues when the URL already has a protocol (e.g., https://example.com)
% resulting in malformed URLs like https://https://example.com.
% This patch checks if the URL already contains "://" and uses it directly if so.
% Note: We use \str_set:Nx (x-expansion) to expand macros like \@homepage first.
\ExplSyntaxOn
\str_new:N \g_yamlresume_url_str

\RenewDocumentCommand{\httpslink}{O{}m}{
  \str_set:Nx \g_yamlresume_url_str {#2}
  \str_if_in:NnTF \g_yamlresume_url_str {://}
    { \url{#2} }
    { \url{https://#2} }
}

\RenewDocumentCommand{\httplink}{O{}m}{
  \str_set:Nx \g_yamlresume_url_str {#2}
  \str_if_in:NnTF \g_yamlresume_url_str {://}
    { \url{#2} }
    { \url{http://#2} }
}
\ExplSyntaxOff

\name{Lucas Martins Oliveira}{}
\title{Engenheiro de Segurança Cibernética | OSCP | Cloud Security}
\phone[mobile]{+55 11 98765-4321}
\email{lucas.oliveira@cybersec.dev}
\homepage{https://lucasoliveira.dev}

\address{Av. Paulista, 1578, Bela Vista, São Paulo, SP, Brasil, 01310-200}{}{}

\extrainfo{{\small \faLinkedin}\ \href{https://www.linkedin.com/in/lucas-martins-oliveira}{@lucas-martins-oliveira} {} {} {} • {} {} {} 
{\small \faGithub}\ \href{https://github.com/lucas-moliveira}{@lucas-moliveira}}

\begin{document}

\maketitle

\section{Informações Básicas}

\cvline{}{\begin{itemize}
\item Engenheiro de segurança cibernética com mais de 8 anos de experiência em proteção de ambientes corporativos e cloud.
\item Especialista em segurança ofensiva, resposta a incidentes e arquitetura de segurança em AWS e Azure.
\item Certificado OSCP, CISSP e AWS Security Specialty, com forte atuação em conformidade (ISO 27001, LGPD, PCI DSS).
\item Liderança técnica de equipes e implementação de programas de segurança desde a concepção até a operação.
\end{itemize}}
\section{Formação}

\cventry{mar. de 2014–dez. de 2018}
        {Bacharelado, Engenharia de Computação, Nota: 8.7/10}
        {Universidade de São Paulo}
        {\url{https://www.usp.br}}
        {}
        {\begin{itemize}
\item Formação sólida em fundamentos de computação, redes e segurança da informação.
\item Desenvolvimento de projetos práticos em segurança de aplicações e análise de vulnerabilidades.
\item Participação em grupo de estudos de cibersegurança e competições de CTF.
\end{itemize}
\textbf{Cursos}: Segurança de Sistemas, Criptografia, Redes de Computadores, Sistemas Operacionais, Banco de Dados, Engenharia de Software, Computação em Nuvem, Inteligência Artificial}
\section{Experiência Profissional}

\cventry{jan. de 2022–Atual}
        {Engenheiro de Segurança Cibernética Sênior}
        {Banco Digital Aurora}
        {\url{https://www.bancaurora.com.br}}
        {}
        {\begin{itemize}
\item Liderança técnica na definição de arquitetura de segurança para serviços em cloud (AWS e Azure).
\item Implementação de programa de threat modeling e revisão de segurança de aplicações.
\item Condução de resposta a incidentes críticos e análise forense.
\item Desenvolvimento de automação de controles de segurança com Terraform e Python.
\item Mentoria de engenheiros júniores e disseminação da cultura de segurança.
\end{itemize}
\textbf{Palavras-chave}: Threat Modeling, Cloud Security, SIEM, Incident Response, Terraform}

\cventry{fev. de 2019–dez. de 2021}
        {Engenheiro de Segurança da Informação}
        {Consultoria RiskTech}
        {\url{https://www.risktech.com.br}}
        {}
        {\begin{itemize}
\item Realização de testes de intrusão em aplicações web, APIs e infraestrutura.
\item Avaliação de conformidade com ISO 27001, PCI DSS e LGPD.
\item Implementação de soluções SIEM e correlação de eventos.
\item Elaboração de relatórios técnicos e planos de remediação.
\item Atuação em projetos de resposta a incidentes para clientes dos setores financeiro e saúde.
\end{itemize}
\textbf{Palavras-chave}: Penetration Testing, ISO 27001, SIEM, LGPD, Risk Assessment}

\cventry{ago. de 2017–dez. de 2018}
        {Estagiário de Segurança da Informação}
        {TechSecure Soluções}
        {\url{https://www.techsecure.com.br}}
        {}
        {\begin{itemize}
\item Monitoramento de alertas de segurança e triagem inicial de incidentes.
\item Apoio na configuração de firewalls e sistemas de detecção de intrusão.
\item Documentação de procedimentos operacionais de segurança.
\item Participação em simulações de phishing e conscientização de usuários.
\end{itemize}
\textbf{Palavras-chave}: Firewalls, IDS/IPS, Phishing, Monitoring}
\section{Idiomas}

\cvline{Português}{Nativo ou Bilíngue}
\cvline{Inglês}{Proficiência Profissional Fluente \hfill \textbf{Palavras-chave}: TOEFL 105, Cambridge C1}
\cvline{Espanhol}{Proficiência Profissional Limitada}
\section{Competências}

\cvline{Segurança Ofensiva}{Avançado \hfill \textbf{Palavras-chave}: Penetration Testing, Metasploit, Burp Suite, OSCP, Nmap}
\cvline{Segurança em Cloud}{Avançado \hfill \textbf{Palavras-chave}: AWS Security, Azure Security, IAM, Terraform, Kubernetes Security}
\cvline{Resposta a Incidentes}{Especialista \hfill \textbf{Palavras-chave}: SIEM, Splunk, Forensics, SOAR, Threat Hunting}
\cvline{Redes e Infraestrutura}{Avançado \hfill \textbf{Palavras-chave}: Firewalls, VPN, IDS/IPS, TCP/IP, Zero Trust}
\cvline{Governança e Conformidade}{Intermediário \hfill \textbf{Palavras-chave}: ISO 27001, LGPD, NIST, PCI DSS, Risk Assessment}
\section{Prêmios}

\cventry{nov. de 2023}
        {Banco Digital Aurora}
        {Profissional de Segurança do Ano}
        {}
        {}
        {Reconhecimento pela liderança na implementação do programa de segurança em cloud e redução de 40\% nos incidentes críticos.}
\section{Certificados}

\cventry{jun. de 2021}
        {Offensive Security}
        {OSCP - Offensive Security Certified Professional}
        {\url{https://www.offsec.com/certifications/oscp/}}
        {}
        {}

\cventry{set. de 2022}
        {ISC2}
        {CISSP - Certified Information Systems Security Professional}
        {\url{https://www.isc2.org/certifications/cissp}}
        {}
        {}

\cventry{mar. de 2023}
        {Amazon Web Services}
        {AWS Certified Security - Specialty}
        {\url{https://aws.amazon.com/certification/certified-security-specialty/}}
        {}
        {}

\cventry{fev. de 2020}
        {CompTIA}
        {CompTIA Security+}
        {\url{https://www.comptia.org/certifications/security}}
        {}
        {}
\section{Publicações}

\cventry{out. de 2023}
        {Detecção de Intrusão em Ambientes Cloud-Native}
        {Revista Brasileira de Segurança Cibernética}
        {\url{https://www.revistabsc.com.br/artigos/deteccao-intrusao-cloud-native}}
        {}
        {Artigo sobre abordagens de detecção de intrusão em ambientes de contêineres e Kubernetes, com foco em machine learning e correlação de eventos.}
\section{Projetos}

\cventry{jan. de 2023–dez. de 2023}
        {Plataforma open-source para agregação e correlação de indicadores de ameaça.}
        {ThreatIntel BR}
        {\url{https://github.com/lucas-moliveira/threatintel-br}}
        {}
        {Desenvolvimento de plataforma para coleta de feeds de threat intelligence e correlação com eventos de segurança; integração com SIEM e geração de alertas contextualizados.
\textbf{Palavras-chave}: Threat Intelligence, Python, SIEM, Open Source}

\cventry{abr. de 2022–out. de 2022}
        {Ferramenta de análise forense de logs para ambientes Linux.}
        {LogHunter}
        {\url{https://github.com/lucas-moliveira/loghunter}}
        {}
        {Criação de ferramenta em Python para análise de logs de autenticação e detecção de comportamentos anômalos; redução do tempo de triagem em 30\%.
\textbf{Palavras-chave}: Python, Forensics, Log Analysis, Linux}
\section{Interesses}

\cvline{Cibersegurança}{CTFs, Bug Bounty, Malware Analysis}
\cvline{Esportes}{Corrida, Natação, Ciclismo}
\cvline{Leitura}{Ficção Científica, Biografias}
\section{Voluntariado}

\cventry{mar. de 2020–Atual}
        {Mentor e Palestrante}
        {OWASP São Paulo Chapter}
        {\url{https://owasp.org/www-chapter-sao-paulo/}}
        {}
        {\begin{itemize}
\item Atuação como mentor em programas de capacitação em segurança de aplicações.
\item Palestras sobre OWASP Top 10 e boas práticas de desenvolvimento seguro.
\item Organização de encontros técnicos e workshops para a comunidade.
\end{itemize}}
    
\end{document}